\section{Return and lease methods in one company book}
\label{sec:shared-return-lease-book}

A customer refund and a lease payment draw on the same company's cash. Separate
method examples can each balance and still fail to describe that competition.
The lease example might spend money that the return example has already paid
to customers. Combining their closing reports would discover the inconsistency
too late: the second payment should have failed when it was proposed.
This makes a common book, rather than a common calling convention, the next
necessary part of the implementation.

The actual Henkel disclosures in the preceding section still lack the contracts
needed for an executable opening. The present example therefore declares a
new synthetic company. It uses retained return and lease calculations to test
shared cash, journals and memory, while leaving actual-company reconstruction
as a separate problem.

\subsection{A common caller preserves different capabilities}

The finite \path{ModuleInvocation/1} interface supplies one request form for
four existing references. Typed returns and shared leases propose transitions;
Henkel measurements assess conditional constraints; the sector reference
produces a report. Their outputs have different meanings. A measurement result
cannot become a posting merely because the caller uses the same envelope.

Each request names a registered module and operation, binds the complete
payload, and identifies the inspected implementation and schema bundle.
The caller selects code from its fixed registry. A supplied path or callback
cannot select another function. Typed-return inputs also retain their original
raw-byte hashes and duplicate-key checks. The complete native result survives,
including journals, gross flows, memory changes and refusal reasons.

An outer status of \path{returned} means that the selected module answered.
Its separate \path{native_status} may be proposed, unsupported, invalid or
unresolved. A request to post the sector report, post a measurement assessment,
or commit a lease proposal is refused before native invocation. Verification
recomputes complete transition and report results. The measurement verifier
instead permits another valid solver witness and another runtime while checking
the underlying equations and all fixed information.

Identity binding does not authenticate source labels. In the retained return
reference, changing a well-formed owner label or packet digest can still produce
a proposal; source admission remains false. The lease reference fixes its
original scope and refuses a changed scope digest. The common caller preserves
both behaviors. It cannot manufacture the missing source-to-state judgment.

\subsection{The original states are incompatible}

The return state belongs to synthetic owner P and begins in 2025. The lease
state belongs to a different synthetic owner and begins in 2024. Their schemas
are closed, their legal-book labels differ, and their account registries contain
22 and 26 accounts respectively. The return state stores recovered lots, tax
adjustments and fiscal history. The lease state stores named rights, their
allocations and a premise that all leases are listed.

Even shared account names are insufficient. Cash needs one owner, currency,
date and opening amount; inventory needs a complete allocation to stock; lease
liabilities need their rights and payment history. A field union does not
resolve those meanings and violates both original validators. Renaming the
lease owner does not repair the original scope restriction. The read-only
compatibility reference consequently rejects direct combination, preserves
unknown membership as unresolved, and accepts only its explicit control:
two valid identical snapshots of the same native profile. That control authorizes
no joint transition.

\subsection{One declared opening and explicit memory ownership}

\path{SyntheticReturnLeaseBook/1} supplies a separate synthetic USD legal book.
It retains the 22 return accounts and explicitly adds the right-of-use asset,
lease liability and lease derecognition gains. Debit balances and expenses are
positive; credit balances and income are negative. The six selected lease
posting categories---cash, right-of-use asset, lease liability, finance expense,
depreciation and other gains---have declared meanings in this 25-account
registry. No mapping between unrelated return cost categories is inferred.

The opening supplies cash 500, inventory 60, a recovery asset of 16, a refund
liability of 40, a right-of-use asset of 120 and a lease liability of 150, together
with the surrounding accounts recorded in the complete fixture. Its signed
original-equity balance is explicitly 20. These are synthetic premises chosen
for the example; the program neither infers that opening from an issuer nor
computes an equity plug. The two lease rights have right-of-use amounts 90 and
30 and liabilities 100 and 50. The declared collection contains exactly those
rights.

The book owns one period, date, cash balance, event history and gross-flow
history. Return methods own the return component, recovered lots and settlement
tax adjustments. Lease methods own the lease-right collection. Other memory,
including tax premises and prior fiscal history, must survive each selected
event unchanged. Exact write sets enforce that separation.

For a dated event $k$, let $B_k$ be the vector of signed account balances before
the event and $J_k$ its journal in the same account order, both measured in USD.
Let $F_k$ contain its gross cash flows, with receipts positive and payments
negative. Pure application checks the exact identities
\[
 B_{k+1}=B_k+J_k,\qquad
 \sum_a J_{k,a}=0,\qquad
 J_{k,\mathrm{cash}}=\sum_{f\in F_k}f.
\]
The first identity applies the journal once. The second preserves the signed
trial balance. The third prevents classified cash flows from disagreeing with
cash. Component allocations are checked as well: recovered lots explain their
inventory, and lease rights explain the right-of-use asset and lease liability.
An aggregate balance alone cannot replace those allocations.

Return settlement and whole-lot sale reuse the checked typed prerequisite,
monetary and memory programs. Their memory schema is extended explicitly to
retain all 25 balances and the lease fields. The lease operations reuse the
original component calculation kernel under the new declared context. The
fixed-owner lease wrapper remains unchanged. This is a separately specified
synthetic instantiation, not a lossless migration of the old owners.

\subsection{Follow both methods through the same cash balance}

At 1 January 2024 the lease reduction terminates right B and remeasures right A.
The journal reduces the right-of-use asset by 20, reduces the credit liability
by 40 and recognizes a gain of 20. No cash moves. The remaining right-of-use
asset is 100 and its liability 110.

On 1 February the return method pays customers 40 and handling costs 8,
releases the refund liability of 40 and recovery asset of 16, and recognizes
recovered inventory of 20. The associated cost-of-sales credit is 4 and selling
expense is 8. The two cash payments remain separate operating flows. A recovered
lot records the units, carrying amount, tax base and acquisition date, while
the settlement tax adjustment remains available for a later tax calculation.

On 1 March the entire recovered lot is sold for 30. Its carrying amount of 20
leaves inventory and enters cost of sales. On 1 January 2025 the lease payment
of 121 comprises interest 11 and principal 110 under the supplied annual
timing assumptions. Interest is an operating outflow and principal a financing
outflow under the retained election. A separately supported depreciation
amount of 10 then reduces the right-of-use asset. Neither that amount nor the
discount rate is inferred from aggregate disclosures.

With no other cash events in the example, the closing amount follows directly:
\[
 \mathrm{cash}_{\mathrm{close}}
 =500-\underbrace{(40+8)}_{\text{refund and handling}}
       +\underbrace{30}_{\text{recovered-stock sale}}
       -\underbrace{121}_{\text{lease payment}}
 =361.
\]
The book closes this sequence with right-of-use assets of 90, no remaining
lease liability and inventory of 60. Five event identities and five gross
flows survive. Return events leave the lease memory unchanged; lease events
preserve the return component, lots, tax adjustments and other return memory.

The low-cash case distinguishes genuine shared execution from two independent
examples. Start with cash 130 and omit the recovered-stock sale. The signed original-equity
amount is supplied as 390 in this alternative opening, keeping the declared
trial balance balanced; the program makes no automatic adjustment. The return
settlement leaves 82, so the later lease payment of 121 fails the shared funding
check. The completed lease reduction and return settlement remain; the failed
payment changes nothing. With the same opening cash but without the return
settlement, the lease payment succeeds and leaves 9. Cash competition therefore
changes the admissible sequence even though each method's arithmetic is
unchanged.

\subsection{What replay establishes}

A proposal binds the whole before-state, event and implementation bundle.
Application recomputes the proposal before returning a new in-memory state.
An old proposal cannot apply to a changed cash balance, and a forged journal,
flow, component change or admission flag fails full comparison. Event identities
cannot repeat. An empty event list still requires a valid opening. Unsupported
currency or tax operations are refused; missing return measurements or lease
membership remain unresolved.

All five selected event kinds are also compared with the original native
methods. The original return comparison stays in its 2025 period and original
owner; the new book expressly uses different return dates and ownership.
Journals and activity/amount cash flows match exactly for the same numerical
component inputs. Component and lot memory match apart from the declared lot
acquisition dates. The comparison supports reuse of those effects, not identity
between the old and new companies.

This reference supplies a finite shared-book execution example. It preserves
tax memory but does not close tax or produce complete financial statements for
the new profile. General module substitution, arbitrary histories, group and
foreign-currency methods, actual-source admission and physical database
durability remain separate requirements. The existing database reference is
tied to another engine; its transaction evidence does not transfer automatically.
The full schemas, complete openings, requests, proposals and checked comparisons
are retained under \path{data/accounting-master/} in
\path{module-invocation-20261024},
\path{shared-state-compatibility-20261024} and
\path{shared-book-successor-20261024}.

\subsection{Sell part of a recovered lot without losing its remaining units}
The whole-lot sale above leaves no recovered stock to track. A company that
sells only part of that stock needs a quantity and cost record for what remains.
The existing lot already records both; its carrying amount alone would not
identify how many units were sold. A separately selected method therefore
takes the sold quantity as an explicit input rather than deriving it from cash
proceeds or inventory value.

This method assumes that the remaining units are homogeneous and divisible,
that their measurement basis has not changed, and that proportional allocation
is the selected treatment. The caller also supplies affirmative findings for
completed performance and the retained equality of tax base and carrying
amount. Those findings record explicit premises. Unknown findings remain
unresolved, and false findings place the event outside this method's scope.
A heterogeneous lot or a different tax treatment requires another method.

Immediately before a sale, let $N>0$ be the remaining units, $C\geq0$ their
carrying amount in USD, and $T=C$ their tax base under the selected treatment.
The event supplies sold units $q$, with $0<q\leq N$, and cash proceeds $P\geq0$
in USD. The fraction applies to the stock immediately before this event:
\[
 f=\frac{q}{N},\qquad
 C_{\mathrm{released}}=fC,\qquad
 T_{\mathrm{released}}=fT.
\]
The remaining lot keeps its identity and acquisition date, with units $N-q$,
carrying amount $C-fC$ and tax base $T-fT$. It becomes disposed exactly when
$q=N$. A positive-quantity lot whose carrying amount is zero remains live;
zero cost does not imply that no goods remain.

The journal debits cash by $P$ and cost of sales by $fC$, and credits revenue
by $P$ and inventory by $fC$. Its debit-positive sum is
$P+fC-P-fC=0$. Cash rises by $P$ and inventory falls by $fC$, so net assets and
profit both rise by $P-fC$. In the indirect operating cash calculation, adding
back the inventory reduction to that profit effect gives
$(P-fC)+fC=P$, the same receipt as the direct cash record. The reduction in
the lot's carrying amount equals the inventory posting; every other lot and
the separate inventory balance remain unchanged. This sale assesses no tax
expense and makes no new measurement of inventory.

After the February settlement in the shared-book example, the recovered lot
has two units with carrying amount and tax base of 20. Cash is 452 and total
inventory is 80. Selling one unit for 15 releases 10 of cost, leaves one unit
with carrying amount and tax base of 10, and gives cash 467 and inventory 70.
The prior lease gain of 20, cost-of-sales credit of 4 and handling expense of 8
give recorded profit of 16 before this sale; its margin of 5 raises profit to
21. Operating cash since the opening is $-48+15=-33$. Assets of 827 less
liabilities of 214 equal equity of 613, also opening equity of 592 plus profit
of 21.

If the remaining unit is then sold for another 15 under the same conditions,
the new fraction is one: one unit is sold from one remaining unit. The second
sale releases the remaining 10, leaves cash 482 and total inventory 60, and
marks the lot disposed. The combined account changes equal those of the
original whole-lot sale for 30. The two event identities and cash receipts
remain separate; equal aggregate effects do not make their histories identical.

A proposal retains the entire before-state, event, allocation and resulting
book. Application recomputes the complete proposal. A changed remaining
quantity, altered allocation or balanced but different journal fails that
comparison. Selling more than the available units, selling a disposed lot,
reusing an event identity or backdating a sale is refused without changing
the successful preceding state.

This optional transition has checked quantity, journal and statement identities
and reproduces the original full-sale state when all units are sold. The
existing whole-lot method remains unchanged. Connecting the optional event to
the mixed-book authoring and annual dispatch interfaces is separate work;
these local controls do not establish that those interfaces already support it.
